\section{Materials and Methods}

\subsection{Cell Culture and Experimental Design}

The evaluation dataset comprises two cell lines (iMC~ISOR544C and iMC~MUTR544C) across three culture conditions, acquired in two independent experiments serving as biological replicates (EXP1, $n = 63$ images; EXP2, $n = 63$ images), for a total of 126 baseline ($t = 0$) wound images. The acquisition followed a full factorial design (2~cell lines $\times$ 3~conditions $\times$ 2~independent experiments); both cell lines were present in both experiments (iMC~ISOR544C: 26~images in EXP1, 40~in EXP2; iMC~MUTR544C: 37~in EXP1, 23~in EXP2), so cell line and experiment are crossed rather than nested in this evaluation dataset. The three conditions are an ALK5 inhibitor at $0.1\,\mathrm{mM}$ (a selective TGF-$\beta$ receptor type~I antagonist, at the concentration selected for testing by the generating laboratory), a DMSO vehicle control, and an untreated media reference. These conditions are used solely as a test bed to probe whether the pipeline resolves condition-dependent closure dynamics; this work makes no claims about the underlying pharmacology.


Confluent monolayers were wounded using a standardized mechanical scratch assay. Treatment media were applied immediately after wounding. Time-lapse imaging began at $t = 0$ (baseline, fully open wound) and continued for approximately $24$ hours, with images acquired at ${\sim}1$-hour intervals, yielding $20$--$25$ timepoints per trajectory. Quantification was performed on cropped fields of approximately $900 \times 735$ pixels, corresponding to a pixel size of $1.24\,\mu$m/pixel (field of view ${\approx}1.12 \times 0.91$~mm).

iMCs were seeded into $96$-well plates at a density of $5 \times 10^{4}$ cells per well in media containing either DMSO or the indicated drug conditions. After overnight incubation, wounds were created in the cell monolayers using the Incucyte\textsuperscript{\textregistered} $96$-Well WoundMaker tool. Cells were washed twice with PBS, after which conditioned media were reapplied. Images were acquired immediately after wounding and every hour for up to $24$ hours using the Incucyte\textsuperscript{\textregistered} Live-Cell Analysis System with a $4\times$ objective (yielding the $1.24\,\mu$m/pixel calibration stated above).

\subsection{Automated Wound Segmentation}

The quantification pipeline operates in two stages: a single-frame detector that estimates wound boundaries independently for each image, followed by a temporal integration stage that enforces a monotonic closure constraint across the time series. The resulting boundary estimates are converted into wound masks and quantitative measurements.

\begin{figure}[H]
	\centering
	\includegraphics[width=1.0\linewidth]{diagram_main_figure}
	\caption{%
		\textbf{Automated scratch-wound quantification pipeline with monotonic temporal integration.}
		\textbf{(A)} Experimental workflow. Confluent cell monolayers cultured in 96-well plates are wounded using a standardized mechanical scratch assay and imaged by phase-contrast time-lapse microscopy. Sequential images capture progressive wound closure from baseline to the final timepoint.
		\textbf{(B)} Single-frame wound detection. Each frame is independently segmented using local variance--based texture analysis to distinguish the wound channel from the surrounding cellular monolayer. Candidate wound boundaries are refined through edge detection, RANSAC-based outlier rejection, and Savitzky--Golay smoothing before reconstruction of the wound mask.
		\textbf{(C)} Temporal integration with monotonic closure constraint. Independent frame-wise segmentations produce wound-area trajectories, wound boundaries are modeled independently for each image column using a Kalman-filter framework with a hard monotonic constraint, ensuring that wound area cannot increase over time ($A(t+1) \leq A(t)$). This generates temporally consistent wound trajectories and strictly non-increasing area profiles.
		\textbf{(D)} Quantitative metrics derived from the final segmentation masks, including wound area ($A(t)$), closure fraction ($f(t)$), edge displacement ($\Delta u(t)$ and $\Delta \ell(t)$), and closure speed ($v(t) = -dA(t)/dt$). The fully automated, annotation-free pipeline enables robust extraction of wound-healing dynamics without user-defined threshold tuning while preserving biologically consistent temporal evolution.%
	}
	\label{fig:pipeline-overview}
\end{figure}

\subsubsection{Problem Formulation}

Let $I(x, y, t) \in [0, 1]$ denote the intensity-normalized grayscale image at spatial
position $(x, y)$ and discrete timepoint $t \in \{0, 1, \ldots, T\}$. The image has dimensions $H \times W$ (rows $\times$ columns).
We seek a binary wound mask:
\begin{equation}
	M(x, y, t) = \begin{cases}
		1 & \text{if pixel $(x, y)$ is wound at time $t$} \\
		0 & \text{otherwise}
	\end{cases}
	\label{eq:mask}
\end{equation}

and per-column edge functions:
\begin{align}
	u(x, t) &= \min\{y : M(x, y, t) = 1\} \quad \text{(upper edge)} \label{eq:upper-edge} \\
	\ell(x, t) &= \max\{y : M(x, y, t) = 1\} \quad \text{(lower edge)} \label{eq:lower-edge}
\end{align}

The wound area at time $t$ is:
\begin{equation}
	A(t) = \sum_{x=0}^{W-1} \max\!\bigl(0,\; \ell(x,t) - u(x,t) + 1\bigr)
	\label{eq:area}
\end{equation}

The monotonic closure assumption states that under normal scratch-wound conditions, wound area is non-increasing over time: cells do not retract after migrating into the wound. This assumption holds in standard assays where closure dynamics dominate over transient imaging artefacts; it is encoded as a computational constraint, not a universal biological law:
\begin{equation}
	\boxed{A(t+1) \leq A(t) \quad \forall\; t \geq 0}
	\label{eq:monotonic}
\end{equation}


\subsubsection{Single-Frame Detection}\label{sec_single_frame}

The single-frame detector maps an input image $I(x,y,t)$ to validated edge functions $u_\text{det}(x,t)$ and $\ell_\text{det}(x,t)$ through a seven-stage pipeline that transforms a raw grayscale image into a geometrically consistent wound mask with smooth boundary curves (Figure~\ref{fig:pipeline-example}).

\begin{figure}[H]
	\centering
	\includegraphics[width=0.9\linewidth]{methods_pipeline_example}
	\caption{%
		\textbf{Representative pipeline output on a single frame.}
		Phase-contrast image acquired at $t = 0$ (ALK5i condition, EXP2).
		The wound channel (center, low local variance) is flanked by cell monolayers
		(high local variance). Detected upper and lower boundaries are overlaid as
		colored lines after RANSAC outlier rejection and Savitzky--Golay smoothing.
		The shaded region corresponds to the reconstructed binary wound mask
		(Equation~\eqref{eq:mask-reconstruction}).%
	}
	\label{fig:pipeline-example}
\end{figure}

\textbf{Stage~1: Local Variance Texture Map.}
Cell monolayers exhibit high spatial variance due to cell borders, nuclei, and phase-contrast
halos, while the wound channel is comparatively smooth \citep{Yin2012} (Figure~\ref{fig:pipeline-example}b).
We use local variance as a texture discriminator \citep{Koenig2017,Loizou2016,Ambrossio2022}.
For a square window of side length $w$ (default $w = 15$), the local variance at each pixel
is computed via the uniform-filter identity:
\begin{equation}
	\sigma^2(x, y) = \langle I^2 \rangle_w(x, y) - \bigl[\langle I \rangle_w(x, y)\bigr]^2
	\label{eq:local-variance}
\end{equation}
where $\langle \cdot \rangle_w$ denotes the box-filter mean over a $w \times w$ neighborhood.
The computation is implemented via vectorized array operations for efficiency.

\textbf{Stage~2: Y-Band Detection.}
The wound is identified as a horizontal band spanning nearly the full image width.
We compute the row-averaged variance profile $\bar{V}(y)$, smooth it with a 1-D Gaussian
kernel ($\sigma_y = 10.0$ pixels), locate the global minimum $y_c$, and extract the
contiguous segment of sub-threshold rows (35th percentile) containing $y_c$,
expanded by a margin fraction of 0.15.

\textbf{Stage~3: Column-Wise Edge Extraction.}
Within the Y-band, upper and lower wound boundaries are extracted for each column.
The edge-detection threshold $\tau$ is derived adaptively from the cell regions flanking the
wound: $\tau = 0.5 \cdot \text{median}(\sigma^2_\text{cell})$.
Edge search is fully vectorized using array boolean operations.

\textbf{Stage~4: RANSAC Outlier Rejection.}
We apply an iterative polynomial-fit-and-reject scheme inspired by RANSAC
\citep{Fischler1981}.
A degree-2 polynomial captures genuine wound curvature while rejecting localized spikes.
At each iteration the polynomial is fit to all currently valid columns; columns whose
absolute residual from the fit exceeds the maximum deviation threshold (25~pixels) are
classified as outliers and replaced by the fitted value, and the polynomial is refit on
the updated edge before the next iteration. Edges with fewer than 10 valid columns are
left unmodified. Three iterations are used.

\textbf{Stage~5: Savitzky--Golay Smoothing.}
Residual high-frequency jitter is removed with a Savitzky--Golay filter \citep{Savitzky1964},
preserving edge shape without phase distortion.
Parameters: window length $w_\text{sg} = 51$, polynomial order $p_\text{sg} = 3$.

\textbf{Stage~6: Mask Reconstruction.}
The binary mask is reconstructed from smooth edges using vectorized broadcasting:
\begin{equation}
	M(x, y) = \bigl[u(x) \leq y \leq \ell(x)\bigr] 
	\label{eq:mask-reconstruction}
\end{equation}

\textbf{Stage~7: Geometric Validation.}
Four quality-control criteria are evaluated on each mask (Supplementary Table~\ref{tab:qc-criteria}). A detection passes QC if and only if all four criteria are satisfied.
When the primary pipeline fails, an optional fallback strategy is activated using constrained
region growing from a seed at $y_c$ with morphological cleanup
(hole filling at 500~px, small-object removal at 1000~px, disk closing at radius~5).



\subsubsection{Temporal Constraint: Kalman-Filtered Monotonic Closure}
\label{sec:temporal-constraint}

A naive implementation of the monotonic closure constraint enforces a hard per-column clamp, such that the upper edge $u(x,t)$ can only move inward relative to $u(x,t-1)$ (and symmetrically for the lower edge). While this deterministic rule strictly encodes the biological axiom (Equation~\ref{eq:monotonic}), it introduces an irreversible failure mode: spurious inward detections—arising from debris, illumination artifacts, or transient focus shifts—are permanently propagated to all subsequent frames, with no mechanism for recovery.

To mitigate this limitation, we replace the hard clamp as the primary update mechanism with a per-column Kalman filter. The filter fuses the detector output with a monotonic prior, weighting each according to a data-driven estimate of observation uncertainty. Noisy or low-contrast detection are down-weighted, while consistent high-confidence observations are retained. 

A hard monotonic constraint is subsequently applied as a secondary safety mechanism to guarantee $A(t) \leq A(t-1)$ under all conditions. With this hierarchical design the Kalman filter handles typical artifact-contaminated observations, while the hard constraint provides an unconditional correctness guarantee for residual failure cases. 

The full per-frame processing sequence of the temporal integrator, including detector input, filtering steps, constraint enforcement, and numerical safeguards, is summarized in Supplementary Table~\ref{tab:temporal-sequence}. Four numerical safeguards (Supplementary Table~\ref{tab:temporal-sequence}, Steps 2, 7, 9, 11) prevent failure-mode propagation across frames; their formal statement follows below.


At each timepoint $t$, the single-frame detector (Section~\ref{sec_single_frame}) produces per-column edge estimates $u_\text{det}(x,t)$ and $\ell_\text{det}(x,t)$, corresponding to the upper and lower wound boundaries, respectively. These estimates are treated as observations and are mapped by the temporal integrator to constrained edges $u(x,t)$ and $\ell(x,t)$ that satisfy the monotonic closure constraint. We denote by $P(x,t)$ the per-column uncertainty associated with the temporal estimate.

At the initial timepoint $t = 0$, the detector outputs are accepted without constraint:
\begin{equation}
	u(x,0) = u_\text{det}(x,0), \qquad \ell(x,0) = \ell_\text{det}(x,0)
\end{equation}
These edges define the reference wound geometry for the trajectory. Initial per-column uncertainty $P(x,0)$ is estimated from the spatial residuals between the RANSAC-rejected and Savitzky--Golay smoothed edges at $t = 0$; derivation, fallback values, and numerical safeguards are detailed in Supplementary Table~\ref{tab:pipeline-parameters}.

\textbf{Kalman-filtered constraint ($t > 0$).}
For subsequent timepoints, the temporal integrator updates the constrained edges by combining the detector observations with the previous state under a monotonicity prior. The update proceeds in three steps.

\emph{Prediction.}
The predicted state is the previous constrained edge; predicted uncertainty grows by process
noise $Q$:
\begin{equation}
	\hat{u}^-(x, t) = u(x, t-1), \qquad P^-(x, t) = P(x, t-1) + Q
	\label{eq:kalman-predict}
\end{equation}

Process noise $Q$ (px$^2$) is estimated empirically from QC-validated trajectories
$\mathcal{C}$ as the variance of consecutive frame-to-frame edge differences:
\begin{equation}
	\hat{Q} = \operatorname{Var}\!\Bigl(\bigl\{\,u_\text{det}(x,t{+}1) - u_\text{det}(x,t)\;\big|\; x,\,t \in \mathcal{C}\bigr\}\Bigr)
	\label{eq:Q-estimate}
\end{equation}
$Q$ is estimated once from QC-validated trajectories and cached for reuse; the
conservative default $Q_0 = 400$~px$^2$ is applied only on the first run, before any
cache exists. Estimation therefore draws on previously processed clean trajectories
rather than on the trajectory being filtered, avoiding a circular dependency.

\emph{Spatial pre-smoothing.}
The raw detected edge is smoothed with a 1-D Gaussian kernel ($\sigma_\text{spatial} = 5$
columns) before the Kalman update to suppress column-level artifacts:
\begin{equation}
	\tilde{u}_\text{det}(x, t) = \bigl(G_{\sigma} * u_\text{det}(\cdot,\, t)\bigr)(x)
	\label{eq:spatial-smooth}
\end{equation}

\emph{Per-column observation noise.}
Observation noise $R(x, t)$ is estimated from three data-driven signals:
\begin{equation}
	R(x, t) = R_\text{base}(t) \cdot C^{-2}(x, t) \cdot J(x, t) \cdot m_\text{method}
	\label{eq:R-estimate}
\end{equation}
where $R_\text{base}(t) = \operatorname{Var}_{x}(u_\text{det} - \tilde{u}_\text{det})$
is the frame-level detector noise, $C(x,t)$ is the local wound-vs-cell variance contrast
(low contrast implies high $R$). The jump penalty $J(x,t)$ penalizes large deviations from the prediction:
\begin{equation}
	J(x, t) = 1.0 + \left(\frac{|\tilde{u}_\text{det}(x,t) - u(x,t-1)|}{\text{median}_x(|\tilde{u}_\text{det}(\cdot,t) - u(\cdot,t-1)|)}\right)^2
	\label{eq:jump-factor}
\end{equation}
where the denominator is the frame-level median edge jump across all columns. The method multiplier $m_\text{method} \in \{1.0, 3.0\}$ further penalizes the less reliable fallback detector method.

\emph{Kalman update.}
The gain $K(x,t) \in [0, 1]$ blends prediction and observation:
\begin{align}
	K(x, t) &= \frac{P^-(x, t)}{P^-(x, t) + R(x, t)} \label{eq:kalman-gain} \\
	u_\text{filt}(x, t) &= u(x, t-1) + K(x, t)\bigl[\tilde{u}_\text{det}(x, t) - u(x, t-1)\bigr] \label{eq:kalman-update} \\
	P_\text{upd}(x, t) &= \bigl(1 - K(x, t)\bigr) \cdot P^-(x, t) \label{eq:P-update}
\end{align}
High $R$ (noisy observation) implies small $K$, causing the filter to rely on the monotonic
prior.
Low $R$ (clean, high-contrast frame) implies large $K$, trusting the detector.

\textbf{Hard safety net} (secondary guarantee).
After the Kalman blend, the biological axiom is enforced unconditionally:
\begin{align}
	u(x, t) &= \max\!\bigl(u_\text{filt}(x, t),\; u(x, t-1)\bigr) \label{eq:hard-upper} \\
	\ell(x, t) &= \min\!\bigl(\ell_\text{filt}(x, t),\; \ell(x, t-1)\bigr) \label{eq:hard-lower}
\end{align}
This guarantees $A(t) \leq A(t-1)$ for all $t > 0$ regardless of Kalman output.

The pipeline ultimately produces, for each trajectory, a sequence of wound masks and the corresponding area time series $A(t)$. The Kalman process adds $O(W)$ operations per frame. Total per-frame complexity remains $O(HW)$, dominated by the variance map computation.

 
\textbf{Wound Closure}.
When the normalized wound area satisfies $\frac{A(t)}{A(0)} < \theta_\text{close}$ with $\theta_\text{close} = 0.01$, the wound is considered fully closed. For all subsequent timepoints, the edge positions are held constant ($u(x,t) = u(x,t-1)$ and $\ell(x,t) = \ell(x,t-1)$), preventing spurious numerical fluctuations from sub-pixel noise in near-zero residual areas. The threshold of 1\% was selected as a practical floor below which pixel-quantized edge positions cannot reliably distinguish residual wound from closed tissue at this imaging resolution. Like the monotonic safety net, this closed-edge hold is a one-sided (absorbing) rule: once a trajectory is declared closed it cannot reopen. Because no trajectory in this cohort reached the $\theta_\text{close}$ floor within the observation window (mean closure ${\sim}17\%$ at 20~h), the rule never activated here and therefore did not influence the reported closure-fraction estimates; its effect would become relevant only in datasets that approach full closure.


\subsubsection{Constraint Ablation Study}
\label{sec:ablation-design}

To decompose the individual contributions of the Kalman filter and the hard monotonic safety net, we defined four constraint arms applied to an identical trajectory set:
\textbf{(A) Baseline} --- raw single-frame detector output with no temporal constraint;
\textbf{(B) Hard} --- deterministic per-column max/min clamp only, no Kalman filter;
\textbf{(C) Kalman-only} --- Kalman-filtered edges (Section~\ref{sec:temporal-constraint}) without the hard monotonic safety net; and
\textbf{(D) Kalman+hard} --- the full production pipeline combining the Kalman filter with the hard safety net.

The ablation was evaluated on the 106-trajectory wound-dynamics cohort (ALK5i: $n=15$; DMSO: $n=50$; Media: $n=41$), spanning both cell lines and two independent experiments.
All trajectories were processed through all four arms from the same raw images; the only algorithmic difference was the temporal constraint stage applied after single-frame detection.

Five metrics were computed per trajectory:
(i) \emph{violation rate} --- fraction of consecutive frame pairs with a non-monotonic wound-area increase ($A(t+1) > A(t)$);
(ii) \emph{smoothness} --- mean absolute second difference (MASD) of the wound-area time series (lower is smoother);
(iii) \emph{frame-to-frame area variance} $\sigma^2_\text{FTF}$;
(iv) \emph{final closure fraction} $(A_0 - A_T)/A_0$; and
(v) \emph{area coefficient of variation} (CV~$= \sigma_A/\bar{A}$).
Bootstrap confidence intervals (1000 resamples; seed fixed at 42 for reproducibility) were computed for each metric mean across the 106 trajectories.


\subsection{Wound Quantification}

Once the wound boundaries have been extracted and temporally constrained, quantitative measures of wound healing can be derived from the resulting area trajectory $A(t)$ and edge positions $u(x,t)$ and $\ell(x,t)$. These metrics represent the inward movement of the wound boundaries at each column $x$. 

Fractional wound closure is defined as the proportion of the initial wound area that has closed over time. It ranges from 0 (no closure) to 1 (complete closure), and is given by
\begin{equation}
	f(t) = \frac{A(0) - A(t)}{A(0)}.
	\label{eq:fractional-closure}
\end{equation}


To quantify spatial dynamics, we measure the displacement of the upper and lower wound edges between the initial timepoint $t = 0$ and a target time $t_T$, defining the edge displacement as
\begin{align}
	\Delta u(x) &= u(x, t_T) - u(x, 0), \label{eq:delta-upper} \\
	\Delta \ell(x) &= \ell(x, 0) - \ell(x, t_T). \label{eq:delta-lower}
\end{align}


Wound healing rate is defined by converting the edge displacements of Equations~\eqref{eq:delta-upper}--\eqref{eq:delta-lower} into velocities through normalization by elapsed time. The mean closure speed represents the average inward movement of the wound boundaries and is defined as
\begin{equation}
	v_\text{mean} = \frac{v_\text{upper} + v_\text{lower}}{2},
	\label{eq:closure-speed}
\end{equation}
where each edge velocity is computed as the average displacement per unit time across the set $\mathcal{V}$ of valid columns:
\begin{equation}
	v_\text{edge} = \frac{1}{|\mathcal{V}|} \sum_{x \in \mathcal{V}} \frac{|\Delta_\text{edge}(x)|}{t_T}.
	\label{eq:edge-velocity}
\end{equation}

Averaging over valid columns suppresses localized edge artifacts without discarding the per-column spatial resolution preserved throughout processing.


\subsection{Statistical Analysis}

Pairwise treatment comparisons versus DMSO were performed using Welch's $t$-test
(no equal-variance assumption), stratified by cell line.
Multiple comparison correction was applied using the Bonferroni method across all four
pairwise comparisons within each metric (two non-control conditions, ALK5i and Media,
$\times$ two cell lines).
Effect sizes were quantified as Cohen's $d$, defined as $d = \frac{\bar{x}_\text{treated} - \bar{x}_\text{control}}{s_\text{pooled}}$.

The time-series analysis was performed using a linear mixed-effects model. The response variable $\mathrm{distance}_{ij}$ for trajectory $i$ at timepoint $j$ was modeled as
\begin{equation}
\begin{aligned}
\mathrm{distance}_{ij}
= {}& \mu
+ \beta_1 t_j
+ \beta_2 \,\mathrm{condition}_i
+ \beta_3 \,\mathrm{cell\_line}_i
+ \beta_4 \,(t_j \times \mathrm{condition}_i) \\
& + \beta_5 \,(t_j \times \mathrm{cell\_line}_i)
+ \beta_6 \,(\mathrm{condition}_i \times \mathrm{cell\_line}_i)
+ \beta_7 \,(t_j \times \mathrm{condition}_i \times \mathrm{cell\_line}_i) \\
& + b_{0i}
+ b_{1i} t_j
+ \varepsilon_{ij}.
\end{aligned}
\label{eq:mixed-model}
\end{equation}

Here, $\mu$ denotes the fixed intercept, $\beta_k$ the fixed-effect coefficients, $b_{0i}$ and $b_{1i}$ the random intercept and slope for trajectory $i$, and $\varepsilon_{ij}$ the residual error term. The response $\mathrm{distance}_{ij}$ is the mean inward edge displacement, calibrated to microns ($1.24\,\mu$m/pixel), and the time index $t_j$ is the imaging frame (equivalently hours, given the ${\sim}1$-hour interval), so the time coefficient $\beta_1$ is expressed in microns per hour ($\mu$m/h). The fixed effects follow a full factorial specification ($t \times \mathrm{condition} \times \mathrm{cell\_line}$), so $\beta_6$ and $\beta_7$ capture the condition~$\times$~cell-line and three-way interactions. Models were fit by restricted maximum likelihood using the BFGS optimizer. Inclusion of the per-trajectory random slope $b_{1i}$ was strongly supported by a likelihood-ratio test against a random-intercept-only model ($\chi^2(2) = 1598.5$, $p < 10^{-12}$), confirming that closure rate varies meaningfully across trajectories.

Model assumptions were evaluated by testing residual normality using the Shapiro--Wilk test and homogeneity of variance using Levene's test (see Supplementary Table~\ref{tab:supp-assumption-checks}).
Fixed-effect estimates from linear mixed-effects models are generally robust to moderate residual non-normality when group sizes are sufficient \citep{Schielzeth2020}; the significant non-normality in the Line~2 groups (Supplementary Table~\ref{tab:supp-assumption-checks}) is therefore expected to have minimal impact on the time main-effect estimate. Welch's correction to the cross-sectional $t$-tests independently accommodates unequal variance without assuming normality.

\subsection{Expert Validation}

Detection accuracy was assessed by a single trained observer who reviewed all 126 publication-condition baseline frames (ALK5i, DMSO, and Media) using a custom image annotation application (Supplementary Figures~\ref{fig:review-app}--\ref{fig:review-app-correct}). The reviewer was blinded to the treatment condition assigned to each image; blinding to cell line identity was not applied, as morphological characteristics of each line are distinguishable under phase-contrast imaging. Each frame was independently classified as correctly identified (wound boundary visually consistent with the pipeline output, or absence of wound confirmed for QC-flagged images) or incorrectly identified. Annotation was conducted under strict criteria: any visible discrepancy in boundary alignment, even if limited to a subset of image columns, was classified as incorrect. The resulting binary classifications served as the reference standard for computing detection accuracy and QC module performance metrics.

\subsection{Implementation}
The pipeline was implemented in Python~3.11 using NumPy, OpenCV, SciPy, and scikit-image \citep{vanderWalt2014}. The pipeline is fully deterministic and reproducible. The per-frame computational complexity is $O(HW)$, dominated by the variance map computation. Across all $7{,}486$ processed frames, the primary detector succeeded on $92.2\%$; the remaining $7.8\%$ invoked the region-growing fallback (Section~\ref{sec_single_frame}, Stage~7). The implementation is available at \url{https://github.com/GiorgioRicciardiello/WoundDetectionBounds}.
